\section{Rotational Fourier Defect and Gaussian Characterization}
\label{sec:rotational-fourier-defect-and-gaussian-characterization}

This section characterizes Gaussian coordinate laws by the rotational
rotational Fourier defect.
The centered Fourier feature defining this defect is bounded.
Its definition uses no logarithm of the characteristic function and remains
valid when the characteristic function has zeros.
The conclusion belongs to the same broad lineage as the Kac--Bernstein
theorem and the classical characterization theory of independent linear
forms \cite{kac1939characterization,kagan1973characterization}.
Its hypothesis is different.
The rotational Fourier defect vanishes at every frequency pair.
This leads directly to a global linear ODE.

\subsection{The centered Fourier feature}
\label{sec:centered-fourier-feature}

This subsection requires only $\E X=0$ and $\E|X|<\infty$.
Neither a density nor a second moment is needed.
Write
$\phi(u)=\E e^{iuX}$, and for independent copies $X_1,X_2$ put
\[
 O_\theta=\begin{pmatrix}\cos\theta&-\sin\theta\\
                         \sin\theta&\cos\theta\end{pmatrix},
 \qquad
 W_{\xi,\eta}(x,y)
   =(e^{i\xi x}-\phi(\xi))(e^{i\eta y}-\phi(\eta)).
\]
This fixes the sign convention for rotation.
Centering removes the two marginal responses.
It leaves only the dependence created by rotation.
In particular,
\begin{equation}
 \E W_{\xi,\eta}(X_1,X_2)=0,
 \qquad
 \E|W_{\xi,\eta}(X_1,X_2)|^2
 =(1-|\phi(\xi)|^2)(1-|\phi(\eta)|^2)\le1.
 \label{eq:feature-product-variance}
\end{equation}
Pointwise, $|W_{\xi,\eta}|\le4$, so its variance under an arbitrary
comparison law is at most $16$.  Both bounds are dimension-free and require
no tail assumption.

\subsection{Gaussian characterization by a global linear ODE}
\label{sec:fourier-ode}

\begin{lemma}[Gaussian characterization by the rotational Fourier defect]
\label{lem:fourier-ode}
For the centered Fourier feature defined above,
\begin{equation}
 R(\xi,\eta)
 =\left.\frac{d}{d\theta}
       \E W_{\xi,\eta}(O_\theta(X_1,X_2))\right|_{\theta=0}
   =\eta\phi'(\xi)\phi(\eta)
          -\xi\phi(\xi)\phi'(\eta).
 \label{eq:rotational-fourier-response}
\end{equation}
Moreover,
\begin{equation}
 R\equiv0
 \quad\Longleftrightarrow\quad
 \Law(X)\text{ is a centered Gaussian, possibly degenerate}.
 \label{eq:fourier-gaussian-characterization}
\end{equation}
\end{lemma}
\begin{proof}
The first moment implies $\phi\in C^1(\R)$.
The joint characteristic function after rotation is
\[
 \phi(\xi\cos\theta+\eta\sin\theta)
 \phi(-\xi\sin\theta+\eta\cos\theta).
\]
The marginal characteristic functions are
$\phi(\xi\cos\theta)\phi(-\xi\sin\theta)$ and
$\phi(\eta\sin\theta)\phi(\eta\cos\theta)$.
Their derivatives at zero vanish because $\phi'(0)=i\E X=0$.
Differentiating the joint characteristic function therefore gives
\eqref{eq:rotational-fourier-response}.

If $R$ vanishes everywhere, continuity and $\phi(0)=1$ allow us to
choose $\eta_0\ne0$ with $\phi(\eta_0)\ne0$.
Fix this single frequency.
The defect equation then becomes the global ODE
\[
 \phi'(\xi)=c\xi\phi(\xi),\qquad
 c=\frac{\phi'(\eta_0)}{\eta_0\phi(\eta_0)},\qquad
 \phi(0)=1.
\]
No division by $\phi(\xi)$ has occurred.
Uniqueness for this linear ODE gives $\phi(\xi)=\exp(c\xi^2/2)$ on all of
$\R$, including any putative zero of $\phi$.
The identity $\phi(-\xi)=\overline{\phi(\xi)}$ forces $c\in\R$.
Also $|\phi(\xi)|\le1$ forces $c\le0$.
Write $c=-\sigma^2$.
The Gaussian conclusion follows.
If $\sigma=0$, we get the point mass at zero.
Substitution proves the converse.
\end{proof}

\subsection{The determinant representation and frequency integrals}
\label{sec:fourier-determinant}

Consider the curve
\[
 u\longmapsto\bigl(u\phi(u),\phi'(u)\bigr)\in\mathbb C^2,
\]
described as the weighted Fourier jet.  Its values satisfy
\begin{equation}
 \det\!\begin{pmatrix}
 \xi\phi(\xi)&\phi'(\xi)\\
 \eta\phi(\eta)&\phi'(\eta)
 \end{pmatrix}
 =-R(\xi,\eta).
 \label{eq:fourier-wedge}
\end{equation}
Thus the rotational Fourier defect is the negative determinant of two values
of this curve.
In the Gaussian case \(\phi'(u)=-\sigma^2u\phi(u)\).
So the curve lies in the fixed line spanned by \((1,-\sigma^2)\).
Conversely, vanishing of every such determinant forces precisely the ODE
used above.

For a non-Gaussian law, some two distinct frequencies have nonzero
defect.
By continuity, a compact neighborhood of this pair already detects the law.
Thus a compact frequency set may be fixed after the coordinate law is chosen.
This set is independent of the ambient dimension.
The diagonal \(\xi=\eta\) always has zero defect and cannot serve this
purpose.
One may also integrate the squared defect over all frequencies:
\begin{equation}
 \iint|R(\xi,\eta)|^2w(\xi,\eta)\,d\xi\,d\eta,
 \qquad
 w>0\ \text{a.e.},\quad
 \iint(1+\xi^2+\eta^2)w(\xi,\eta)\,d\xi\,d\eta<\infty.
 \label{eq:fourier-defect-energy}
\end{equation}
Indeed, \(|R(\xi,\eta)|\le\E|X|(|\xi|+|\eta|)\).
Hence this integral is finite and is zero exactly for centered Gaussian laws.
Integrability of $w$ alone need not make the integral finite.

The same feature gives a local Hellinger lower bound for one rotation.
Choose a phase $\alpha$ with
$\operatorname{Re}(e^{-i\alpha}R(\xi,\eta))>0$ and apply
\[
 |\E_\lambda F-\E_\nu F|
 \le h(\lambda,\nu)
             \{2\E_\lambda F^2+2\E_\nu F^2\}^{1/2}
\]
to $F=\operatorname{Re}(e^{-i\alpha}W_{\xi,\eta})$.
It follows that
$h((O_\theta)_\#\Law(X)^{\otimes2},
\Law(X)^{\otimes2})\ge c|\theta|$ for small
$|\theta|$.
A first moment suffices for this lower bound.
The dimension-free matrix estimates below require the additional
second-moment and Sobolev hypotheses.

\subsection{Dimension-free variance and affine normalization}
\label{sec:variance-and-affine-normalization}

Equation \eqref{eq:feature-product-variance} is the variance estimate used
for empirical or finite-frequency detection.
Averaging \(N\) independent pairs gives standard error at most \(N^{-1/2}\)
under the reference product.
The rough bound \(16/N\) remains available under an arbitrary
comparison distribution.
In particular, the centered Fourier feature does not require a fourth moment,
symmetry, or a density.

\paragraph{Centering, scaling, and phase.}
Let $Y$ have mean \(m\) and characteristic function \(\phi_Y\), and put
\(X=(Y-m)/\sigma\), where \(\sigma>0\).
For \(a=\xi/\sigma\) and \(b=\eta/\sigma\), the rotational Fourier defect
of $X$ satisfies
\begin{equation}
 \begin{aligned}
 R(\xi,\eta)
 &=e^{-i(m/\sigma)(\xi+\eta)}
 \bigl[b\phi_Y'(a)\phi_Y(b)-a\phi_Y(a)\phi_Y'(b)\\
 &\hspace{8em}+im(a-b)\phi_Y(a)\phi_Y(b)\bigr].
 \end{aligned}
 \label{eq:centering-phase}
\end{equation}
The expression in brackets is the derivative at zero of the expectation of
the centered Fourier feature after rotating two independent copies of $Y$.
Its last term comes from differentiating the marginal terms.
Thus centering rotates this mean-corrected expression by
\(-(m/\sigma)(\xi+\eta)\).
Scaling also changes the frequencies.
At fixed corresponding frequencies its modulus is unchanged.
Omitting the marginal correction would invalidate this phase identity.
For a centered symmetric law the defect is real.
At a particular frequency pair, an asymmetric law can still have a real or
zero defect.
So the phase is not, by itself, a complete measure of skewness.

For the concrete comparison $Y\sim\operatorname{Gamma}(2,1)$,
$m=2$, $\sigma=\sqrt2$, and $(\xi,\eta)=(1,2)$, the expression in brackets in
\eqref{eq:centering-phase} equals $4/27$, so the standardized law has
\[
 R(1,2)=\frac4{27}e^{-3i\sqrt2}.
\]
At this particular pair, the entire phase of the rotational Fourier defect is
therefore the centering phase, namely $-3\sqrt2$ modulo $2\pi$.
The variance-one Laplace law has rotational Fourier defect $R(1,2)=-4/27$.
Their moduli agree at these frequencies, while their principal phases
are approximately $116.915^\circ$ and $180^\circ$, respectively.
